\documentclass[answers]{otagoexam}
% Class options:
%   answers - shows answers (remove option to hide answers)
%   embargoed - embargoed exam mode
%   manualpagebreaks - disables automatic page breaks between questions  

% Extra packages
\usepackage{graphicx}


%These fields are required
\department{Department}
\papercode{Paper code}
\papertitle{Paper title}
\paperyear{2023}
\semester{Teaching period} %Semester One/Semester Two/Summer School
\timelimit{3 HOURS}


%If you don't define any of these following, defaults show here will be used 

% Instructions: layout your instructions here.  You can use \showmarks{<marks>} command to show marks without counting them towards the total.  The \totalmarks{} command gives a sum of all marks from \addmarks[<mark>] and \items[<mark>]commands.  If you'd want auto-summed marks, place your manual mark value in the command below - ex. \totalmarks{60} to show 60 marks regardless of the auto-sum.

\instructions{
  Candidates should answer \textbf{All} questions

  \vspace{0.3cm}
  Marks are shown thus: \showmarks{2}

  \vspace{0.3cm}
  The total number of marks for this exam is \totalmarks{}
}

% Material provided (default)
\providedmaterial{Nil}

% Calculators (default)
\calculators{\nocalculators} 

% Other options for calculators
%\calculators{\anymodel} %Any model without networking
%\calculators{\listA}    %List A calculators
%\calculators{\listB}    %List B calculators

% Permission for additional material (default)
\permittedmaterial{No additional material}

% Other instructions (used only in non-embargoed mode)
\otherinstructions{Nil}

% Embargo instructions (used only in embargo mode)
\embargoinstructions{Hand in entire exam attached to answer book(s)}





\begin{document}

% Title page
\maketitle


% Uncomment this if you need physical constants page (can reference it with \ref{physicalconstants})
%\physicalconstants

% Uncomment this if you need formulae page (can reference it with \ref{formulae})
%\formulae

% Section without a title, marked.
\section{}

To add a section to the exam use the \verb$\section{}$ command.  Sections are numbered with upper-case letters.  Exams scripts do not have to have sections.  To specify marks for the whole section use \verb$\addmarks{<mark>}$ command at the end of the section paragraph.\addmarks{12}

\begin{question}

    Create a single question within the \texttt{\textbackslash begin\{question\}}...\texttt{\textbackslash end\{question\}} environment.  Questions are numbered with Arabic numerals.

    \begin{subquestions}
        \item List sub-questions using the \texttt{subquestions} environment within the question text.
        \item Sub-questions are numbered with lower-case letters.

        \item{
            \begin{subquestions}
                \item Nesting \texttt{subquestion} environments is fine.
                \item Second level sub-questions are numbered with Roman numerals.
            \end{subquestions}
        }
        
    \end{subquestions}


\end{question}



% Section with a title
\section{Title}

To create a section with a title use the \verb$\section{<title text>}$ command.

% Marked question
\begin{question}

    Question numbering continues regardless of the sections.  To specify marks for the entire question use the \texttt{\textbackslash addmarks\{<mark>\}} command at the end of the question paragraph.\addmarks{6}

    \begin{subquestions}
        \item List sub-questions using the \texttt{subquestions} environment environment within the question text.
        \item Sub-questions are numbered with lower-case letters.

        \item{
            \begin{subquestions}
                \item Nesting \texttt{subquestion} environments is fine.
                \item Second level sub-questions are numbered with Roman numerals.
            \end{subquestions}
        }
        
    \end{subquestions}
\end{question}

% Marked sub-questions
\begin{question}
    
    Each question starts on a new page.  To create a question with marks per sub-question specify marks per item of the \texttt{subquestions} environment using the \texttt{\textbackslash item{[}mark{]} <item text>} syntax.

    \begin{subquestions}
        \item[3] List sub-questions using the \texttt{subquestions} environment within the question text.
        \item[2]  Sub-questions are numbered with lower-case letters.

        \item{
            \begin{subquestions}
                \item[1] Nesting \texttt{subquestion} environments is fine.
                \item[4] Second level sub-questions are numbered with Roman numerals.   
            \end{subquestions}
        }
        
    \end{subquestions}

\end{question}

% Question with answers
\begin{question}

    To include answers use the \texttt{\textbackslash answer\{<answer text>\} command} within a question or under a sub-question.  Answers are printed in red and they show up on the pdf only when \texttt{answers} option is passed to the \texttt{otagoexam} document class.

    \begin{subquestions}
        \item[3] Sub-question

        \answer{You can provide answer for each sub-question.
            \begin{itemize}
                \item To include a mark break-down for an answer, use the \texttt{\textbackslash showMarks\{<marks>\} command} \showmarks{1.5} 
                \item The marks in \texttt{\textbackslash showMarks\{<marks>\} command} do not affect the marks total. \showmarks{1.5} 
            \end{itemize}
        }

        \item[2] Another sub-question
    \end{subquestions}

    \answer{To add an answer at the end of the question, make sure its within the question body.}

    \answer{All answers are removed in the (no \texttt{answers} option) \texttt{otagoexam} document class mode.}

\end{question}

% This command adds the "please do not turn over" page at the end 
% inserting a blank page (if needed) so that the "please do not turn over"
% message always lands on an even page (back of the exam).
\turnoverpage


\end{document}
